\documentclass[11pt]{article}
\usepackage[a4paper,margin=27mm]{geometry}
\usepackage{amsmath,amssymb,amsthm,xcolor,booktabs,array}
\newcommand{\PROVED}{\textcolor{green!40!black}{\textbf{PROVED}}}
\newcommand{\OPEN}{\textcolor{red!70!black}{\textbf{OPEN}}}
\newcommand{\AUDIT}{\textcolor{orange!70!black}{\textbf{AUDIT}}}
\newtheorem{theorem}{Theorem}
\newtheorem{proposition}{Proposition}
\newtheorem{lemma}{Lemma}
\newtheorem{definition}{Definition}
\title{The universal Witt--Hodge envelope of the odd Poisson object (v98)}
\date{13 August 2026}

\begin{document}
\maketitle

\section{Purpose and type correction}

Put
\[
 E:=\mathcal H_-^0=\mathcal H_-/Z\mathcal H_\cap
\]
and let $\pi_0$ be its centrally normalized scaling representation.  The
goal of this note is to construct a positive Hilbert object from the Witt
adjoint without using $H_T^{-1}$, zeros of $\xi$, or the sign of the Weil
form.

There is a type distinction which is essential for this construction.
Row (b) supplies the injective algebraic map
\[
 J:\bigoplus_{r\ge1}\mathbb C\phi_r\longrightarrow\mathcal C_{\mathbb R},
 \qquad J\phi_r=\delta_r,
\]
and row (c) sends $\delta_r$ to the multiplier $U_r$ on a test-function
space.  The completed odd object $E$, however, is introduced later as a
closed quotient in the category of nuclear Fr\'echet representations.
The existing compatibility theorem compares composition and the scalar
contact $\Lambda(r)$; it does not define a map
\begin{equation}
       W_{\rm alg}\longrightarrow E.                 \tag{1}
\end{equation}
Thus (1), and not merely a choice of norm on $W_{\rm alg}$, has to be
constructed if the Witt metric is to complete $E$.

\section{A source-defined universal construction}

Let $G=\mathbb Q_+^\times$.  The half-Tate normalization of v96 gives the
involution $q^*=q^{-1}$ after replacing $U_q$ by the centered operator
$S_q=q^{1/2}U_q$.  Algebraically $S$ is a representation of $G$ on $E$.

Choose any continuous Hilbertian seminorm $p$ on the nuclear Fr\'echet
space $E$ and let $H_p$ be the Hilbert completion of $E/\ker p$.  Form
\[
 \mathfrak F_p:=\ell^2(G)\widehat\otimes H_p.
\]
On its algebraic core let $\mathfrak R_p$ be the span of the covariance
relations
\begin{equation}
 \delta_{gq}\otimes x-\delta_q\otimes S_gx,
 \qquad g,q\in G,\ x\in E,                           \tag{2}
\end{equation}
whenever the second tensor is interpreted in $H_p$.  Put
\begin{equation}
 \mathscr K_p:=\mathfrak F_p/\overline{\mathfrak R_p},
 \qquad Q_p x=[\delta_1\otimes x].                   \tag{3}
\end{equation}

\begin{theorem}[Free Witt--Hodge quotient --- \PROVED]
The space $\mathscr K_p$ is a positive Hilbert space.  Left translation on
$\ell^2(G)$ descends to a unitary representation $\mathsf S_p$ on
$\mathscr K_p$, and
\begin{equation}
       Q_pS_g=\mathsf S_p(g)Q_p.                     \tag{4}
\end{equation}
Consequently the integrated test action satisfies the Witt--Rosati
identity
\begin{equation}
       \mathsf S_p(f^\vee)=\mathsf S_p(f)^*.         \tag{5}
\end{equation}
The null space
\begin{equation}
 \mathfrak N_p:=\ker Q_p
 =\{x:\delta_1\otimes x\in\overline{\mathfrak R_p}\} \tag{6}
\end{equation}
is precisely the part of $E$ lost when the algebraic covariance relations
are Hilbert-completed.
\end{theorem}

\begin{proof}
The closed quotient in (3) is Hilbert.  Left translation $L_g$ is unitary
on $\ell^2(G)$ and carries every generator of (2) to another generator;
the same is true for $L_{g^{-1}}$.  Hence $\overline{\mathfrak R_p}$ is
reducing and $L_g$ descends unitarily.  Relation (2) with $q=1$ is exactly
(4).  Integrating the unitary group representation gives (5).  Formula
(6) is the definition of the kernel of the quotient map and records the
only possible loss in passing from algebraic to Hilbert covariance.
\end{proof}

Taking a countable separating family $(p_j)$ of Hilbertian seminorms and
the Hilbert direct sum of suitably rescaled $\mathscr K_{p_j}$ produces a
canonical source-defined map
\begin{equation}
 Q_{\rm WH}:E\longrightarrow\mathscr K_{\rm WH},
 \qquad
 \ker Q_{\rm WH}=\bigcap_j\mathfrak N_{p_j}.         \tag{7}
\end{equation}
This is the universal Witt--Hodge envelope associated with the supplied
Fr\'echet topology.  It is constructed before any spectral description.

\section{What Hilbert completion does to a nonunitary mode}

The construction (3) permits an exact calculation of its radical.

\begin{theorem}[Hyperbolic modes are forced into the radical --- \PROVED]
Suppose $0\ne x\in E$ and, for some $g\in G$ of infinite order,
\begin{equation}
                 S_gx=\lambda x.                    \tag{8}
\end{equation}
If $|\lambda|\ne1$, then $Q_px=0$ for every $p$ for which $p(x)\ne0$.
More generally, no positive Hilbert completion which is injective on $x$
and satisfies the adjoint identity $S_g^*=S_{g^{-1}}$ can contain (8)
unless $|\lambda|=1$.
\end{theorem}

\begin{proof}
Restrict the first factor in $\mathfrak F_p$ to the copy of
$\ell^2(\mathbb Z)$ generated by $g$.  Relation (2) says that the class of
$(L_g-\lambda)\delta_1\otimes x$ is zero.  If $|\lambda|\ne1$, the bounded
operator $L_g-\lambda I$ is invertible because the spectrum of the
bilateral shift is the unit circle.  Therefore
$\delta_1\otimes x$ belongs to the closed relation space and $Q_px=0$.

For the second assertion, in any positive Hilbert completion the adjoint
identity makes $S_g$ unitary.  Taking norms in (8) gives
$\|x\|=\|S_gx\|=|\lambda|\|x\|$; injectivity on $x$ and $x\ne0$ imply
$|\lambda|=1$.
\end{proof}

For the centered scaling character attached by row (c) to a zero $\rho$,
$\lambda=g^{\rho-1/2}$.  The theorem therefore kills exactly every
off-critical character.  This conclusion was not inserted in the
definition of $Q_{\rm WH}$; it is forced by positivity and the Witt
adjoint.

There is a second, subtler loss.  The regular module $\ell^2(G)$ has no
eigenvectors, so the closed range of $L_g-\lambda$ is dense also when
$|\lambda|=1$.  Thus the particular free regular quotient (3) can collapse
unitary point characters as well.  To retain them one must take the direct
sum of all positive cyclic Witt quotients, equivalently the universal
$C^*(G)$ representation.  This refinement still kills (and must kill)
all $|\lambda|\ne1$ modes.

\section{The maximal positive Witt quotient}

Let $\mathscr P(E)$ be the set of continuous positive semidefinite
sesquilinear forms $b$ on $E$ satisfying
\begin{equation}
 b(S_gx,y)=b(x,S_{g^{-1}}y)\qquad(g\in G).           \tag{9}
\end{equation}
Define
\begin{equation}
 \mathfrak R_{\rm Hod}(E):=igcap_{b\in\mathscr P(E)}\operatorname{rad}b.
                                                               \tag{10}
\end{equation}

\begin{theorem}[Maximal positive Witt--Hodge separated quotient --- \PROVED]
The quotient
\begin{equation}
                 E_{\rm Hod}:=E/\mathfrak R_{\rm Hod}(E)       \tag{11}
\end{equation}
is separated by its family of positive GNS Hilbert maps.  It is maximal in
the following sense: every continuous map from $E$ to a positive Hilbert
space which intertwines the Witt adjoint factors through (11).  If a
countable separating subfamily $(b_j)$ can be rescaled so that
$\sum_j b_j(x,x)<\infty$ continuously on $E_{\rm Hod}$, its diagonal GNS
map is a faithful embedding into one Hilbert direct sum.  Every
off-critical scaling eigenvector belongs to
$\mathfrak R_{\rm Hod}(E)$.
\end{theorem}

\begin{proof}
Quotient each form $b$ by its radical and complete.  Equation (9) makes
the centered $G$-action unitary on that completion.  The family of all
resulting maps has common kernel equal to the intersection (10), hence
separates (11).  Any other positive intertwining map pulls its Hilbert
inner product back to a member of $\mathscr P(E)$, so its kernel contains
(10), proving maximality.  Under the summability hypothesis stated in the
theorem, the diagonal map has squared norm $\sum_jb_j(x,x)$ and is the
claimed single Hilbert embedding.  The last assertion follows from the
preceding theorem, applied to every $b$.
\end{proof}

This is a genuine universal family of Hilbert completions produced from the
Witt adjoint and the odd source representation, not from $H_T^{-1}$ or from
a spectral positive part.  Producing one global Hilbert completion requires
the countable summability clause in the theorem.  The remaining issues are
that clause, injectivity of $E\to E_{\rm Hod}$, and preservation of the
nuclear character.

\section{The exact separation theorem required for row (d)}

Write $\tau_-(h)=\operatorname{Tr}_E\pi_-(h)$ for the odd nuclear character
of row (c).  On the test convolution algebra set
\begin{equation}
              q_-(f,g):=\tau_-(g^\vee*f).            \tag{12}
\end{equation}

\begin{theorem}[Hilbert descent/Weil positivity equivalence --- \PROVED]
The following statements are equivalent on the primitive test ideal.
\begin{enumerate}
\item $q_-(f,f)\ge0$ for every primitive $f$.
\item The primitive odd character is of positive type and hence has its
      canonical GNS Hilbert realization with
      $\pi_-(f^\vee)=\pi_-(f)^*$.
\item The row-(c) supertrace satisfies
      $B_{\rm nuc}(f,f)\le0$ for every primitive $f$.
\end{enumerate}
When the GNS realization also realizes the nuclear trace as an operator
trace, one has
\begin{equation}
 B_{\rm nuc}(f,f)
 =-\operatorname{Tr}\bigl(\pi_-(f)\pi_-(f)^*\bigr)\le0.        \tag{13}
\end{equation}
\end{theorem}

\begin{proof}
Statements (1) and (2) are the GNS theorem applied to (12).  For primitive
$f$ the two even Tate matrix coefficients vanish, so the supertrace is
minus the odd trace:
\[
 B_{\rm nuc}(f,f)=-\tau_-(f*f^\vee)=-q_-(f,f).
\]
This proves equivalence with (3).  If the integrated operator is nuclear
in the Hilbert realization, the adjoint identity gives (13).
\end{proof}

\begin{corollary}[Exact new target --- \PROVED]
Rows (b) and (c) close row (d) if one proves the following source-labelled
separation statement:
\begin{equation}
 \boxed{
 \mathfrak R_{\rm Hod}(\mathcal H_-^0)=0
 \quad\text{and the induced Hilbert character equals }\tau_- .}           \tag{14}
\end{equation}
No additional Gamma normalization or finite-window sign computation is
needed after (14).
\end{corollary}

\begin{proof}
Here (14) includes the countably summable separating family required in the
preceding theorem.  It then gives a faithful positive Hilbert completion of
the odd object with the Witt adjoint.  Character
compatibility and vanishing of the even Tate part on the primitive ideal
give (13), hence row (d).  Conversely, any faithful Hilbert completion of
the required kind supplies separating forms in $\mathscr P(E)$ and the
same character, so it implies (14).
\end{proof}

\section{Audit of what rows (b) and (c) currently prove}

The contact maps already present in the manuscript do not prove (14).
They give
\[
 W_{\rm alg}\xrightarrow{J}\mathcal C_{\mathbb R}
 \xrightarrow{\mathfrak j}\mathcal D'(\mathbb R_+^\times)
 \longrightarrow\mathsf{Mult}(\mathcal S_>)
\]
and compare the scalar determinant contact with $\Lambda(n)$.  The odd
quotient $E$ is then introduced independently as a nuclear Fr\'echet
quotient.  There is no arrow from the first line onto $E$, and scalar
contact equality cannot imply separation of vectors in $E$.

Accordingly, the following ledger is forced.

\begin{center}
\begin{tabular}{>{\raggedright\arraybackslash}p{.53\linewidth}
                >{\centering\arraybackslash}p{.15\linewidth}
                >{\raggedright\arraybackslash}p{.22\linewidth}}
\toprule
Statement & Status & Location \\
\midrule
Witt--Rosati involution before completion & \PROVED & v96 \\
Free positive Witt--Hodge quotient & \PROVED & (2)--(7) \\
Maximal separated family of positive Witt quotients & \PROVED & (9)--(11) \\
Assembly into one faithful Hilbert completion & \OPEN & summable separating family \\
Off-critical modes lie in every positive radical & \PROVED & Theorem 2 \\
Rows (b)/(c) define a point-separating map into that quotient & \OPEN & (14) \\
Preservation of the odd nuclear character & \OPEN & (14) \\
Row (d), condition $G$ & \OPEN & immediate after (14) \\
\bottomrule
\end{tabular}
\end{center}

The construction therefore produces the missing Hilbert envelope and
locates the only admissible non-circular theorem: prove, directly from the
Poisson--Witt source maps, that their positive invariant forms separate the
odd quotient and that their Hilbert character is Meyer's nuclear character.
Declaring this separation from the scalar identities $\Lambda(n)$ would
assume the required promotion and is not a proof.

\end{document}
